\section{The problem}
\label{sec:problem}

A researcher wants to know whether something worked. A language policy, a curriculum reform, a public campaign, an editorial standard. The intervention has a date, usage before and after it can be counted in a corpus, and the counts have moved. The difference looks like the effect.

It rarely is, and the reason is a measurement problem before it's an estimation problem. A corpus frequency is a score, and like any score it stands in for a construct it doesn't equal. Validity theory names two ways a score can fail the construct it's meant to represent: \term{construct underrepresentation}, where the measure captures less than was intended, and \term{construct-irrelevant variance}, where the measure moves for reasons that have nothing to do with the construct \citep{messick_1989_validity,messick1995Validity}. A corpus frequency carries both. The second is what breaks before-and-after designs. The sources a corpus draws on shift over time, so the construct-irrelevant part of the score doesn't sit still: it moves systematically from one period to the next, in the same direction, and for reasons no reader of the counts can see. A two-point comparison can't separate that movement from the effect the researcher is after.

Corpora are built in a way that guarantees the movement. A corpus isn't a sample of a population's usage; it's a sample of texts, gathered by editors and archivists whose priorities shift \citep{biber1993}, so balance and representativeness are approximations rather than guarantees \citep{McEneryHardie2012}. Diachronic work compounds it, because a matched corpus has to hold its sampling frame steady while the genres inside it drift \citep{LeechHundtMairSmith2009}. So a corpus frequency carries two things at once, what the target population did and what the corpus contributed, and the DiD machinery has to be told which of them it's being asked to identify.

The causal reading enters easily all the same, often through a verb, usage that \mention{responds to} a reform rather than an argued claim. Researchers working on prescription have long felt this pull: confronting normative dictates with usage, it tends to find the causal link assumed rather than demonstrated, from relativizer choice in edited English prose \citep{HinrichsSzmrecsanyiBohmann2015} to five centuries of French injunction set against the record of speech \citep{PoplackDion2009}. One robustness check is grammatical. Delete the causal verb and see what claim is left.

A rise after an intervention is a correlation in time. What licenses a causal reading is a design, and the standard design for this situation is difference-in-differences (DiD): set the treated series against a comparison series the intervention left untouched, and read the effect off the change in their gap. That reading rests on one assumption, called \term{parallel trends}: absent the intervention, the two series would have moved together, so the comparison series stands in for what the treated series would have done.

Because the outcome is a score rather than the construct, the assumption has to hold twice. The rate at which people use the form has to move in parallel across the two series, and so does whatever the corpus itself contributes to the score. Section~\ref{sec:estimand} works this through in numbers: the first can hold exactly while the second fails, and the failure alone produces a four-point effect where the true effect is zero. Causal inference from text is an active field \citep{keith2020text,grimmer2022text,egami2022texts}, but it has mostly handled text as a treatment, outcome, or confounder in a cross-section. Less settled is what DiD's assumptions become when the outcome is a corpus frequency measured across time.

The contribution is a discipline for claiming with DiD, not a tutorial on running it; the econometric side already has a current practitioner's guide of its own \citep{BakerEtAl2026}. A before-and-after corpus contrast supports a causal reading only after three things have been pulled apart: what population the claim is about, what the corpus's composition contributed, and what the measurement procedure contributed. A second separation, among the units at which language varies, at which treatment varies, and at which observations stop being independent, is set out in Section~\ref{sec:estimand}. Conflate any two and the apparent effect can be an artifact of the corpus, not a fact about the language.

Language research has spent a decade on reproducibility, replicability, and robustness, corpus work included \citep{SchweinbergerHaugh2025,Flanagan2025}. That programme asks whether a descriptive estimate survives reanalysis, resampling, and analytic choice \citep{EgbertBiberGrayLarsson2025,Sonning2024}. The question here sits in front of it: whether a before-and-after contrast licenses a causal reading at all. A result can be perfectly reproducible and still identify nothing.

Seven questions drive the paper, and the sections take them in turn. What's the shock, and at what unit does treatment vary (Sections~\ref{sec:estimand} and~\ref{sec:treatments})? Which population rate is the target, and what outcome measures it (Sections~\ref{sec:estimand} and~\ref{sec:outcomes})? What comparison makes the counterfactual plausible, what corpus-specific threat could counterfeit the effect, and what diagnostic would expose it (Sections~\ref{sec:threats} and~\ref{sec:diagnostics})?

The running case is the feminization of profession nouns under francophone language policy \citep{Fujimura2005,BurnettBonami2019}. It's a teaching vehicle, not a finding: the aim is the method, and the case is built so that its verdict is sometimes that the effect can't be identified. A corpus DiD is worth most when it disciplines a causal claim downward, and that's the argument the rest of the paper makes.
